% ===========================================================================
%  离散数学与结构（I） 数论初步 课堂笔记
%  整除与偏序 · gcd 与扩展欧几里得 · 算术基本定理 · 同余与剩余类
%  欧拉定理 · 中国剩余定理 · 欧拉函数求和恒等式
%  Compile:  latexmk -xelatex main      (or: xelatex main, twice)
% ===========================================================================
\documentclass[cjk]{loom}

\newcommand{\N}{\mathbb{N}}
\newcommand{\Z}{\mathbb{Z}}
\newcommand{\Ns}{\mathbb{N}^{*}}            % 正整数
\DeclareMathOperator{\lcm}{lcm}
\newcommand{\cop}{\perp}                     % 互素记号 a \cop b
\newcommand{\gcdp}{\mathop{\mathrm{gcd'}}\nolimits}  % gcd'（撇号不漂浮）

\runningthread{离散数学与结构 · 数论初步}

\begin{document}

\loomcover
  {数论初步：整除、同余与中国剩余定理}
  {离散数学与结构（I）· 课堂笔记}
  {}
  {2026 年 9 月}

\begin{strand}
本讲内容为：整除关系及其偏序结构；最大公因数的三种刻画（最大公因子、最小正线性
组合、扩展欧几里得算法的输出）与 Bézout 恒等式；素数与算术基本定理；同余与
模运算、剩余类与既约剩余系；欧拉定理与费马小定理；中国剩余定理；欧拉函数
的求和恒等式 $\sum_{d\mid m}\varphi(d)=m$；以及欧拉定理的应用——RSA 加密算法。
\end{strand}

\clearpage
{\small\tableofcontents}

\clearpage
% ===========================================================================
\section{整除与偏序}\label{sec:div}

本节的论域为 $\N$ 或 $\Z$（会逐处注明），只使用四则运算 $+,-,\times,\div$ 的
语言。记 $\Ns=\N\setminus\{0\}=\{1,2,3,\ldots\}$ 为正整数集。

\begin{definition}[整除]
对整数 $a,b$，定义
\[
  a\mid b \quad\Longleftrightarrow\quad \exists\, c\in\Z,\ b=ac.
\]
此时称 $a$ 是 $b$ 的\keyword{因子}，$b$ 是 $a$ 的\keyword{倍数}。
\end{definition}

\block{整除是一个偏序（论域为 $\N$ 时）}
\begin{itemize}[leftmargin=1.3em,itemsep=1pt]
  \item \emph{自反性}：$a\mid a$（取 $c=1$）；
  \item \emph{传递性}：$a\mid b\ \wedge\ b\mid c \Rightarrow a\mid c$；
  \item \emph{反对称性}：论域为 $\N$ 时，
        $a\mid b\ \wedge\ b\mid a \Rightarrow a=b$；
        论域为 $\Z$ 时只能得到 $a=\pm b$（例如 $2\mid-2$ 且 $-2\mid2$）。
\end{itemize}
因此整除关系在 $\N$ 上构成\keyword{偏序关系}。在这个偏序下，$1$ 是“最小”的
（$1$ 整除一切），$0$ 是“最大”的（一切整除 $0$）——这与通常的大小关系相反，
注意体会“偏序”与“数值大小”是两回事。我们不去考虑 $0\mid 0$ 是否成立之类的
corner case。

\begin{definition}[最大公因数与最小公倍数]
对整数 $a,b$（不全为 $0$），定义
\[
  \gcd(a,b)=\max\{c\in\N : c\mid a\ \wedge\ c\mid b\},
  \qquad
  \lcm(a,b)=\min\{c\in\Ns : a\mid c\ \wedge\ b\mid c\}.
\]
\end{definition}

两式中的集合都非空且有界，故最值存在：公因子集合含 $1$，且以 $\max(|a|,|b|)$
为上界；正公倍数集合含 $|ab|$。

\block{最大公因数的另一种刻画}
定义
\[
  \gcdp(a,b)=\min\{\,am+bn : m,n\in\Z,\ am+bn>0\,\},
\]
即 $a,b$ 的\keyword{最小正整系数线性组合}。下一节将证明
$\gcdp(a,b)=\gcd(a,b)$——这正是 Bézout 恒等式的内容，而证明的工具是
扩展欧几里得算法。

% ===========================================================================
\section{最大公因数与扩展欧几里得算法}\label{sec:gcd}

本节证明上节预告的等式 $\gcd(a,b)=\gcdp(a,b)$。证明分三步，分别给出三个量
之间的不等式，最后夹出相等。记 $\mathrm{GCD}(a,b)$ 为下面算法的输出。

\block{第一步：$\gcd(a,b)\le\gcdp(a,b)$}
令 $d=\gcd(a,b)$，并设 $\gcdp(a,b)=am'+bn'$。因为 $d\mid a$ 且 $d\mid b$，
显见 $d\mid am'+bn'$，即 $\gcd(a,b)\mid\gcdp(a,b)$，于是
\begin{equation}
  \gcd(a,b)\le\gcdp(a,b). \tag{1}
\end{equation}

\block{第二步：扩展欧几里得算法}
\textbf{算法 $\mathrm{GCD}$（扩展欧几里得算法）。}
输入 $0\le a\le b$；输出 $(d,m,n)$，满足 $d=am+bn$。
\begin{itemize}[leftmargin=1.6em,itemsep=1pt]
  \item 若 $a=0$，输出 $(b,\,0,\,1)$（此时 $d=b=0\cdot a+1\cdot b$）；
  \item 否则做带余除法 $b=aq+r$（$0\le r<a$），递归计算
        $(d,m,n)=\mathrm{GCD}(r,a)$，输出 $(d,\;n-qm,\;m)$。
\end{itemize}
递归的正确性由
\[
  d=rm+an=(b-aq)m+an=a(n-qm)+bm
\]
保证；第一参数从 $a$ 严格降为 $r$，故算法终止。由于输出 $d$ 是 $a,b$ 的
正整系数线性组合，由 $\gcdp$ 的最小性得
\begin{equation}
  \gcdp(a,b)\le\mathrm{GCD}(a,b). \tag{2}
\end{equation}

\block{第三步：$\mathrm{GCD}(a,b)\le\gcd(a,b)$}
对递归深度归纳可证 $\mathrm{GCD}(a,b)\mid a$ 且 $\mathrm{GCD}(a,b)\mid b$：
基例 $a=0$ 时输出 $b$，显然整除 $0$ 与 $b$；递归步中 $d\mid r$、$d\mid a$，
而 $b=aq+r$，故 $d\mid b$。于是 $\mathrm{GCD}(a,b)$ 是 $a,b$ 的公因子，
不超过最大公因子：
\begin{equation}
  \mathrm{GCD}(a,b)\le\gcd(a,b). \tag{3}
\end{equation}

\begin{theorem}[Bézout 恒等式]\label{thm:bezout}
设 $a,b$ 不全为 $0$，则
\[
  \gcd(a,b)=\gcdp(a,b)=\min\{am+bn>0:m,n\in\Z\}.
\]
特别地，存在整数 $m,n$ 使 $\gcd(a,b)=am+bn$，且这样的系数可由扩展欧几里得
算法实际算出。
\end{theorem}

\begin{proof}
结合 (1)(2)(3) 得
$\gcdp(a,b)\le\mathrm{GCD}(a,b)\le\gcd(a,b)\le\gcdp(a,b)$，故三者相等。
\end{proof}

\begin{corollary}[互素判据]\label{cor:coprime-crit}
$\gcd(a,b)=1$ 当且仅当存在整数 $m,n$ 使 $am+bn=1$。
\end{corollary}

\begin{definition}[互素]
若 $\gcd(a,b)=1$，称 $a$ 与 $b$ \keyword{互素}（coprime），记作 $a\cop b$。
\end{definition}

\begin{lemma}[互素的乘积封闭性]\label{lem:cop-product}
若 $\gcd(a,b)=1$ 且 $\gcd(a,c)=1$，则 $\gcd(a,bc)=1$。
\end{lemma}

\begin{proof}
由推论 \ref{cor:coprime-crit}，存在 $m,n,m',n'\in\Z$ 使
\[
  am+bn=1,\qquad am'+cn'=1.
\]
两式相乘（写成 $bn=1-am$、$cn'=1-am'$ 的形式）得
\[
  bc\cdot nn'=(1-am)(1-am')=1-a(m+m'-amm'),
\]
即
\[
  a\,(m+m'-amm')+bc\,(nn')=1.
\]
再用一次推论 \ref{cor:coprime-crit}，得 $\gcd(a,bc)=1$。
\end{proof}

\begin{remark}
这条引理是中国剩余定理唯一性论证的关键：由它归纳可得，若
$m_1,\ldots,m_t$ 两两互素且每个 $m_i$ 都整除 $x$，则乘积
$m_1\cdots m_t$ 也整除 $x$。
\end{remark}

% ===========================================================================
\section{素数与算术基本定理}\label{sec:fta}

\begin{definition}[素数]
正整数 $p$ 是\keyword{素数}（prime），当且仅当
\[
  p>1\ \wedge\ \bigl(a\mid p \Rightarrow a=1\ \vee\ a=p\bigr),
\]
即 $p$ 的正因子只有 $1$ 与 $p$ 自身。
\end{definition}

\begin{theorem}[算术基本定理]\label{thm:fta}
任何 $a\in\Ns$，都存在唯一的（非降）素数序列
$p_1\le p_2\le\cdots\le p_n$，使
\[
  a=p_1p_2\cdots p_n.
\]
（$a=1$ 时理解为空乘积，定义为 $1$。）
\end{theorem}

\begin{proof}
\textbf{存在性}（强归纳）。$a=1$ 是空乘积。设 $a>1$：若 $a$ 本身是素数，
取单一项序列即可；否则 $a=bc$ 且 $1<b,c<a$，对 $b,c$ 分别用归纳假设，
把两个素数分解合并重排即得 $a$ 的分解。

\textbf{唯一性}（对 $a$ 作归纳）。设
\[
  a=p_1p_2\cdots p_n=q_1q_2\cdots q_m,
\]
两个序列均按非降排列。若 $p_1=q_1$，约去这个公共因子得到更小的数，
由归纳假设两个分解完全相同。

否则，不妨设 $p_1<q_1$。两者是\emph{不同的}素数，故 $\gcd(p_1,q_1)=1$；
由 Bézout 恒等式，存在 $s,t\in\Z$ 使
\[
  p_1s+q_1t=1.
\]
两边同乘 $q_2q_3\cdots q_m$（注意 $q_1q_2\cdots q_m=a$）：
\[
  s\,p_1q_2\cdots q_m \;+\; ta \;=\; q_2q_3\cdots q_m.
\]
左边第一项显然被 $p_1$ 整除；第二项 $ta$ 也被 $p_1$ 整除（因为 $p_1\mid a$）。
所以 $p_1\mid q_2q_3\cdots q_m$。但 $q_2q_3\cdots q_m<a$，由归纳假设它的素因子
分解唯一，其素因子只有 $q_2,\ldots,q_m$，全都 $\ge q_1>p_1$——矛盾。
\end{proof}

\begin{remark}
唯一性证明中“$p_1\mid q_2\cdots q_m$ 导致矛盾”这一步，本质上就是
Euclid 引理（素数整除乘积则整除某个因子）的归纳形式；而 Euclid 引理的
标准证明同样以 Bézout 恒等式为核心。
\end{remark}

% ===========================================================================
\section{同余与模运算}\label{sec:cong}

以下固定模数 $m$，并在不发生混淆时省略 “$\mathrm{mod}\ m$”。

\begin{definition}[同余]
\[
  a\equiv b\pmod m \quad\Longleftrightarrow\quad m\mid a-b.
\]
\end{definition}

\block{基本性质}
同余是等价关系，并且与加法、乘法相容：
\begin{itemize}[leftmargin=1.3em,itemsep=1pt]
  \item $a\equiv a$；
  \item $a\equiv b \Rightarrow b\equiv a$；
  \item $a\equiv b\ \wedge\ b\equiv c \Rightarrow a\equiv c$；
  \item $a\equiv b\ \wedge\ c\equiv d \Rightarrow a+c\equiv b+d$，
        且 $ac\equiv bd$。
\end{itemize}

\block{与普通四则运算的区别：消去律不一定成立}
模运算中不能随意约去两边的公因子。例如 $2\cdot1\equiv2\cdot4\pmod 6$，
但 $1\not\equiv4\pmod 6$。什么时候可以消去？答案是：当这个因子\emph{可逆}时。

\begin{theorem}[逆元的存在性]\label{thm:inverse}
若 $\gcd(c,m)=1$，则存在 $c^{-1}$ 使
\[
  c\cdot c^{-1}\equiv1\pmod m.
\]
\end{theorem}

\begin{proof}
由 Bézout 恒等式（推论 \ref{cor:coprime-crit}），存在 $c^{-1},t\in\Z$ 使
$cc^{-1}+mt=1$，即 $cc^{-1}\equiv1\pmod m$。
\end{proof}

\begin{corollary}[消去律]
若 $\gcd(c,m)=1$，则
\[
  ac\equiv bc\pmod m \;\Longrightarrow\; a\equiv b\pmod m.
\]
\end{corollary}

\begin{proof}
两侧同乘 $c^{-1}$ 即可。
\end{proof}

% ===========================================================================
\section{剩余类、欧拉函数与欧拉定理}\label{sec:euler}

\block{等价类与代表元}
模 $m$ 的同余关系把 $\Z$ 分成等价类：
\[
  \overline{0}=\{\ldots,-m,0,m,\ldots\},\quad
  \overline{1}=\{\ldots,-m+1,1,m+1,\ldots\},\quad\ldots\quad
  \overline{i}=\{i+km : k\in\Z\}.
\]
取每一类的最小非负代表元，得到\keyword{完全剩余系}
\[
  \Z_m=\{\overline{0},\overline{1},\ldots,\overline{m-1}\}.
\]
在剩余类上规定运算
\[
  \overline{a}+\overline{b}=\overline{a+b},\qquad
  \overline{a}\cdot\overline{b}=\overline{ab}.
\]
由同余与加乘的相容性（第 \ref{sec:cong} 节），运算结果不依赖代表元的选取，
即运算是良定义的。

\begin{definition}[既约剩余系与欧拉函数]
\[
  \Z_m^{*}=\{\,\overline{c}\in\Z_m : \gcd(c,m)=1\,\}
\]
称为模 $m$ 的\keyword{既约剩余系}；其中每个元素都可逆，且
$\overline{a}^{-1}=\overline{a^{-1}}$。其元素个数
\[
  \varphi(m)=|\Z_m^{*}|
\]
称为\keyword{欧拉函数}。
\end{definition}

\begin{theorem}[欧拉定理]\label{thm:euler}
对任意 $n,a$ 满足 $\gcd(n,a)=1$（即 $a\cop n$），有
\[
  a^{\varphi(n)}\equiv1\pmod n.
\]
\end{theorem}

\begin{proof}
因为 $a$ 在模 $n$ 下可逆，映射 $\overline{i}\mapsto\overline{i}\cdot\overline{a}$
是 $\Z_n^{*}$ 上的一个双射（逆映射为乘以 $\overline{a}^{-1}$），即“乘以 $a$”
只是把既约剩余系中的元素重新排列。于是
\[
  \prod_{i\in\Z_n^{*}} i \;\equiv\; \prod_{i\in\Z_n^{*}}(i\cdot a)
  \;\equiv\; a^{\varphi(n)}\prod_{i\in\Z_n^{*}} i \pmod n.
\]
乘积 $\prod i$ 仍属于 $\Z_n^{*}$，故可逆；两侧同乘其逆即得
\[
  1\equiv a^{\varphi(n)}\pmod n. \qedhere
\]
\end{proof}

\begin{corollary}[费马小定理]\label{cor:fermat}
对素数 $p$，若 $p\nmid a$，则
\[
  a^{p-1}\equiv1\pmod p.
\]
\end{corollary}

\begin{proof}
$p$ 为素数时 $\varphi(p)=p-1$（$1,\ldots,p-1$ 都与 $p$ 互素），
代入欧拉定理即得。
\end{proof}

\begin{corollary}\label{cor:fermat-all}
对素数 $p$ 与任意整数 $a$，
\[
  a^{p}\equiv a\pmod p.
\]
\end{corollary}

\begin{proof}
若 $p\nmid a$，在费马小定理两边同乘 $a$；若 $p\mid a$，两边都 $\equiv0$。
\end{proof}

% ===========================================================================
\section{中国剩余定理}\label{sec:crt}

\begin{theorem}[中国剩余定理，CRT]\label{thm:crt}
设正整数 $m_1,m_2,\ldots,m_t$ 两两互素（$i\ne j\Rightarrow m_i\cop m_j$），
记 $m=m_1m_2\cdots m_t$。则对任意
$(a_1,a_2,\ldots,a_t)\in\Z_{m_1}\times\Z_{m_2}\times\cdots\times\Z_{m_t}$，
存在唯一的 $a\in\Z_m$，使
\[
  a\equiv a_i\pmod{m_i}\qquad(i=1,2,\ldots,t).
\]
\end{theorem}

\begin{proof}
记 $M_i=m/m_i$（$i=1,\ldots,t$）。由于 $m_1,\ldots,m_t$ 两两互素，
$\gcd(M_i,m_i)=1$，故 $M_i$ 在模 $m_i$ 下有逆，记为 $M_i'$，即
$M_iM_i'\equiv1\pmod{m_i}$（定理 \ref{thm:inverse}）。

\textbf{存在性}（构造）。我们验证
\[
  x\equiv\sum_{i=1}^{t}a_iM_iM_i'\pmod m
\]
是解。实际上，当 $i\ne j$ 时 $m_j\mid M_i$，于是和式中除第 $j$ 项外全部
被 $m_j$ 整除：
\[
  \sum_{i=1}^{t}a_iM_iM_i'\equiv a_jM_jM_j'\equiv a_j\pmod{m_j}.
\]

\textbf{唯一性}。假设 $x_1,x_2$ 均为方程组的解，则
$x_1\equiv x_2\pmod{m_i}$（$1\le i\le t$），即每个 $m_i$ 都整除 $x_1-x_2$。
因为 $m_1,\ldots,m_t$ 两两互素，由引理 \ref{lem:cop-product} 归纳可得
乘积 $m$ 整除 $x_1-x_2$，即 $x_1\equiv x_2\pmod m$。
\end{proof}

\begin{example}[《孙子算经》——物不知数]
今有物不知其数，三三数之剩二，五五数之剩三，七七数之剩二，问物几何？
即求解
\[
  x\equiv2\pmod3,\qquad x\equiv3\pmod5,\qquad x\equiv2\pmod7.
\]
按 CRT 证明的构造：$m=3\cdot5\cdot7=105$，分别找出
$35$ 的倍数中模 $3$ 余 $1$ 的、$21$ 的倍数中模 $5$ 余 $1$ 的、
$15$ 的倍数中模 $7$ 余 $1$ 的数，比如 $70,21,15$，则所有解为
\[
  x\equiv2\times70+3\times21+2\times15=233\equiv23\pmod{105}.
\]
\end{example}

\begin{remark}
CRT 也可以表述为一个环同构：
\[
  \Z_m\;\longrightarrow\;\Z_{m_1}\times\cdots\times\Z_{m_t},\qquad
  \overline{x}\longmapsto
  \bigl(\overline{x},\ldots,\overline{x}\bigr),
\]
即“模乘积的剩余类”与“各模数剩余类的元组”一一对应，且保持加法与乘法。
可逆元恰好对应各坐标均可逆的元组，因此
$\Z_m^{*}\cong\Z_{m_1}^{*}\times\cdots\times\Z_{m_t}^{*}$，
特别地 $|\Z_m^{*}|=\prod_i|\Z_{m_i}^{*}|$，即欧拉函数是积性的：
$a\cop b\Rightarrow\varphi(ab)=\varphi(a)\varphi(b)$。
\end{remark}

% ===========================================================================
\section{欧拉函数的求和恒等式}\label{sec:phisum}

\begin{theorem}[Gauss 恒等式]\label{thm:gauss-sum}
对任意正整数 $m$，
\[
  \sum_{d\mid m}\varphi(d)=m.
\]
\end{theorem}

\begin{proof}
把 $\Z_m$ 中的元素按它与 $m$ 的最大公因数分组（各组两两不交）：
\[
  m=|\Z_m|
  =\sum_{d\mid m}\bigl|\{i\in\Z_m : \gcd(i,m)=d\}\bigr|.
\]
对固定的 $d\mid m$，映射 $i\mapsto i/d$ 给出双射
\[
  \{i\in\Z_m : \gcd(i,m)=d\}\;\longleftrightarrow\;\Z_{m/d}^{*},
\]
因为 $\gcd(i,m)=d$ 当且仅当 $i=dj$ 且 $\gcd(j,m/d)=1$。故这一组恰有
$\varphi(m/d)$ 个元素，于是
\[
  m=\sum_{d\mid m}\bigl|\Z_{m/d}^{*}\bigr|
   =\sum_{d\mid m}\varphi\Bigl(\frac{m}{d}\Bigr)
   =\sum_{d\mid m}\varphi(d),
\]
最后一步只是对求和指标做代换 $d\leftrightarrow m/d$（$d$ 跑遍 $m$ 的所有
因子时，$m/d$ 也跑遍所有因子）。
\end{proof}

\begin{remark}[欧拉函数的计算（补充）]
对素数幂，模 $p^k$ 下与 $p^k$ 不互素的数恰为 $p$ 的倍数，故
\[
  \varphi(p^k)=p^k-p^{k-1}\qquad(k\ge1).
\]
结合 CRT 给出的积性（第 \ref{sec:crt} 节末的注记），对
$n=\prod_{i=1}^{r}p_i^{k_i}$ 有
\[
  \varphi(n)=\prod_{i=1}^{r}\bigl(p_i^{k_i}-p_i^{k_i-1}\bigr)
  =n\prod_{i=1}^{r}\Bigl(1-\frac{1}{p_i}\Bigr).
\]
例如 $\varphi(36)=36\bigl(1-\tfrac12\bigr)\bigl(1-\tfrac13\bigr)=12$。
\end{remark}

% ===========================================================================
\section{RSA：欧拉定理的应用}\label{sec:rsa}

仿射密码、置换密码等经典密码都是\keyword{对称加密}：加密与解密用
同一把密钥，因此通信前必须设法秘密地交换密钥。RSA（Rivest--Shamir--Adleman，
1977）是\keyword{非对称加密}：加密用公开的公钥，解密用自己保存的私钥，
私钥全程不必传送。

\block{密钥生成}
\begin{itemize}[leftmargin=1.3em,itemsep=1pt]
  \item 取两个很大的不同素数 $p,q$，计算 $N=pq$；
  \item $\varphi(N)=(p-1)(q-1)$，选取满足 $\gcd\bigl(e,\varphi(N)\bigr)=1$ 的
        正整数 $e$（$1\le e<\varphi(N)$），$(e,N)$ 作为\keyword{公钥}公开；
  \item 由扩展欧几里得算法求 $d$ 使 $de\equiv1\pmod{\varphi(N)}$，
        $(d,N)$ 作为\keyword{私钥}自己保存。
\end{itemize}

\block{加密与解密}
明文与密文都表示为 $0\le x<N$ 的整数：
\[
  y\equiv x^{e}\pmod N\quad(\text{加密}),\qquad
  x\equiv y^{d}\pmod N\quad(\text{解密}).
\]

\begin{theorem}[RSA 的正确性]
沿用上面的记号，对任意 $0\le x<N$，都有
\[
  x^{ed}\equiv x\pmod N,
\]
即解密变换是加密变换的逆变换。
\end{theorem}

\begin{proof}
由 $de\equiv1\pmod{\varphi(N)}$，存在正整数 $t$ 使 $de=t\varphi(N)+1$。

若 $\gcd(x,N)=1$，由欧拉定理（定理 \ref{thm:euler}）$x^{\varphi(N)}\equiv1
\pmod N$，于是
\[
  x^{ed}=x^{t\varphi(N)+1}=\bigl(x^{\varphi(N)}\bigr)^{t}\cdot x\equiv x\pmod N.
\]

若 $\gcd(x,N)>1$（且 $x\ne0$），则 $\gcd(x,N)=p$ 或 $q$，不妨设为 $p$。
此时 $x$ 与 $q$ 互素，由费马小定理（推论 \ref{cor:fermat}）
$x^{q-1}\equiv1\pmod q$，从而
\[
  x^{ed-1}=\bigl(x^{q-1}\bigr)^{t(p-1)}\equiv1\pmod q,
  \quad\text{即}\quad x^{ed}\equiv x\pmod q;
\]
而 $p\mid x$，两边模 $p$ 都是 $0$，即 $x^{ed}\equiv x\pmod p$。
由 CRT（或 $p,q$ 为不同素数）得 $x^{ed}\equiv x\pmod{pq}$。
$x=0$ 时结论显然。
\end{proof}

\begin{example}[一个可手算的例子]
取 $p=5,q=11$，则 $N=55$，$\varphi(N)=40$。取 $e=3$，由
$3\cdot27=81\equiv1\pmod{40}$ 得 $d=27$。
\begin{itemize}[leftmargin=1.3em,itemsep=1pt]
  \item 加密 $x=7$：$y\equiv7^{3}=343\equiv13\pmod{55}$；
        解密 $13^{27}\equiv7\pmod{55}$；
  \item 非单位明文 $x=5$：$y\equiv5^{3}=125\equiv15\pmod{55}$；
        解密 $15^{27}\equiv5\pmod{55}$。
\end{itemize}
\end{example}

\begin{remark}[安全性与数字签名]
要解密就需要 $d$，而求 $d$ 通常需要先知道 $\varphi(N)$，也就是要分解
$N$；目前尚未发现分解大整数的快速算法，这是 RSA 安全性的基础（但
“破解 RSA”与“整数分解”的等价性并未被证明）。另外，加密与解密可以交换
次序（$e_K(d_K(x))=x$），把 RSA 倒过来用——用私钥“加密”、大家用公钥
“解密”验证——就得到\keyword{数字签名}。这里的裸 RSA 只保留数论原理，
实际使用还需要编码与填充方案。
\end{remark}


\end{document}
